%قبل از شروع مسئله ترجیحاً چیزی ننویسید. اما اگر خواستید بنویسید، در قالب لیستی (فقرات و شمارش) بنویسید.
%برای اضافه کردن نمره، بعدی از شروع‌مسئله، [x] را اضافه کنید و به جای x  نمره مورد نظر را قرار دهید. برای زیرمسئله‌ها هم به همین صورت است.
%\شروع‌مسئله[10]
\شروع‌مسئله [20]
\متن‌سیاه{خوب، بد، زشت}

دکتر شکیبا از دوران کودکی حروف کوچک انگلیسی را به دو دستۀ خوب و بد تقسیم می‌کند. او یک رشته از حروف را زیبا می‌نامد، اگر حداکثر $k$ حرف بد در آن یافت شود.
به ما یک رشته از حروف کوچک انگلیسی داده‌شده است. قصد داریم تعداد زیررشته‌های زیبای متمایز آن را مطابق با استانداردهای دکتر شکیبا پیدا کنیم.

\متن‌سیاه{نکات:}
\شروع{فقرات}
\فقره
در این سوال مجاز به استفاده از توابع \متن‌لاتین{hash} پایتون نیستید.
\فقره
در این سوال اجازه استفاده از اپراتور $==$ و ابزارهای مشابه برای مقایسه رشته‌ها را ندارید.
\پایان{فقرات}
\متن‌سیاه{ورودی:}\\
اولین خط ورودی شامل رشته‌ی $s$ است. 
در خط بعدی یک دنباله‌ی ۲۶تایی از $0$ و $1$ به شما داده می‌شود که خوب یا بد بودن حروف الفبا را نشان می‌دهد. (خوبی با $1$ و بدی با $0$ نشان داده می‌شود)
سپس در خط آخر عدد $k$ به شما داده می‌شود. 

$$1 \leq |s| \leq 1500$$
$$1 \leq k \leq 1500$$

\متن‌سیاه{خروجی:}\\
در تنها خط خروجی باید تعداد زیررشته‌های زیبای متمایز رشته‌ی $s$ را چاپ کنید.

\متن‌سیاه{ورودی و خروجی نمونه:}\\
\begin{latin}
	\begin{tabularx}{\textwidth}{|X | X|}
		\hline
		\textbf{input} &
		\textbf{output}\\
		\hline
		\texttt{ababab}
		&
		\texttt{5}\\
		\texttt{01000000000000000000000000}&\\
		\texttt{1}&\\
		\hline
		
	\end{tabularx}
\end{latin}


%اگر مسئله چند قسمت داشت به این صورت عمل کنید:
%\شروع{زیرمسائل}
%\زیرمسئله[0] % مسئلۀ ۱آ
%\زیرمسئله % مسئلۀ ۱ب
%\زیرمسئله % مسئلۀ ۱پ
%\زیرمسئله % مسئلۀ ۱ت
%\پایان{زیرمسائل}

\solution{
در هر زیر رشته باید در $O(1)$ بفهمیم که این زیر رشته زیباست یا نه. برای این کار ابتدا یک  پیش‌پردازش از $O(n)$ باید روی تمام  زیر رشته ها انجام بدهیم و با استفاده از  \متن‌لاتین{partial sum} تعداد حروف بد را از ابتدا تا مکان $x$ حساب کنیم. سپس برای هر رشته کافی است عدد انتهای رشته را از عدد ابتدای رشته کم کنیم و جواب را به‌دست آوریم حال که می‌توانیم برای هر زیر رشته چک کنیم که زیباست یا نه. در انتها کافی است روی تمام زیر رشته‌ها با $O(n^2)$ حرکت کنیم و اگر خوب بود درهم‌سازی آن را به مجموعه زیر رشته‌های خوب اضافه کنیم (چون خود آن می‌تواند بسیار بزرگ باشد و در حافظۀ ما جای نگیرد) همچنین درهم‌سازی هر زیر رشته را می‌توان در $O(1)$ حساب کرد. کافی است هر کاراکتر که اضافه می‌شود، عدد مربوط به آن را اضافه کنیم و هر کاراکتر که کم می‌شود عدد مربوط به آن را از مقدار درهم‌سازی فعلی کم کنیم و در مواقعی که کاراکتر از سمت چپ کم می‌شود علاوه بر کم کردن درهم‌سازی آن باید مقدار کل درهم‌سازی را تقسیم بر یک عدد نیز بکنیم.

کد آقای وحید زهتاب با نام \متن‌لاتین{GBU\_sol.py} به عنوان نمونه در پوشه مربوطه قرار داده شده. 

}



%%%%%%%%%%%%%%%%%%%%%%%%%%%%%%%%%%%%%%%%%%%%%%%%%%%%%
% برخی دستورات مفید برای تمرین‌ها. %
%%%%%%%%%%%%%%%%%%%%%%%%%%%%%%%%%%%%%%%%%%%%%%%%%%%%%

% فرمول‌های نمونه:
%$f(x) = \Theta(x)$
%$T(x, y) \leq \log(x) + 2^y + \binom{2n}{n}$
% {\tt code\_function}


% لیست موارد بدون شماره‌گذاری:
%\شروع{فقرات}
%    \فقره مورد اول 
%        \شروع{فقرات}
%            \فقره مورد اول داخلی 
%                \شروع{فقرات}
%                    \فقره اولین مورد خیلی داخلی 
%                    \فقره دومین مورد خیلی داخلی 
%                \پایان{فقرات}
%            \فقره مورد دوم داخلی 
%        \پایان{فقرات}
%    \فقره مورد دوم 
%\پایان{فقرات}

% لیست موارد شماره‌گذاری شده:
%\شروع{شمارش}
%    \فقره یک مورد 
%    \فقره یک مورد دیگر 
%    \فقره باز هم یک مورد دیگر
%\پایان{شمارش}

% برای نوشتن یک فرمول چند خطی که یک شماره بخوره. برای حالت‌های دیگر به سایت زیر مراجعه کنید:
% https://www.overleaf.com/learn/latex/Aligning_equations_with_amsmath
%\begin{align} 
%    \begin{split}
%        (x+y)^3 &= (x+y)^2(x+y) \\
%                &= (x^2+2xy+y^2)(x+y) \\
%                &= (x^3+2x^2y+xy^2) + (x^2y+2xy^2+y^3) \\
%                &= x^3+3x^2y+3xy^2+y^3
%    \end{split}                                 
%\end{align}

% نوشتن ماتریس و آرایه دوبعدی
%\begin{align}
%    A = 
%    \begin{bmatrix}
%        A_{11} & A_{21} \\
%        A_{21} & A_{22}
%    \end{bmatrix}
%\end{align}

% اضافه کردن تصویر با شرح. برای اضافه کردن تصویر بدون شرح، فقط دستور درج‌تصویر را استفاده کنید.
%\شروع{شکل}[h!]
%\تنظیم‌ازوسط
%  \درج‌تصویر[پهنا=0.5\پهنای‌متن]{img.jpg}
%\شرح{شرح تصویر}
%\برچسب{تصویر:۱}
%\پایان{شکل}
%\پایان{زیرمسائل}

%کد زیر را ببینید.
%%\lstinputlisting{source_filename.py}
%\begin{latin}
%\begin{lstlisting}
%# Import the modules
%import sys
%import random
%
%ans = True
%
%while ans:
%question = raw_input("Ask the magic 8 ball a question: (press enter to quit) ")
%
%answers = random.randint(1,8)
%
%if question == "":
%sys.exit()
%
%elif answers == 1:
%print "It is certain"
%
%elif answers == 2:
%print "Outlook good"
%
%elif answers == 3:
%print "You may rely on it"
%
%elif answers == 4:
%print "Ask again later"
%
%elif answers == 5:
%print "Concentrate and ask again"
%
%elif answers == 6:
%print "Reply hazy, try again"
%
%elif answers == 7:
%print "My reply is no"
%
%elif answers == 8:
%print "My sources say no"
%\end{lstlisting}
%\end{latin}
